% -------------------------------------------------------------------------
% WBS ID: RES-ML-INT
% Deliverable: Integration with Simulation and Optimisation
% Status: In progress
% Evidence status: Integration deferred by current hydrodynamic scope
% Review status: Not reviewed
% Last updated: 2026-07-19
% Dependencies: RES-ML-SUR, RES-VER-SPEC, RES-OPT-SPEC
% Downstream interfaces: Software implementation planning, RES-ML-SPEC
% -------------------------------------------------------------------------

\section{RES-ML-INT --- Integration with Simulation and Optimisation}
\label{sec:res-ml-int}

\subsection{Purpose and Scope}
\label{sec:res-ml-int-purpose}

This Deliverable records a deferred integration path. During the
active project phase, Software integration is limited to data
handling, structured outputs and validation reports for the regional
and local hydrodynamic models. Surrogate, optimisation and FSI
integration remain future work.

\subsection{Research Questions}
\label{sec:res-ml-int-research-questions}

% TODO: State the questions that must be answered before the Deliverable
% can be accepted.

\subsection{Terminology and Definitions}
\label{sec:res-ml-int-definitions}

% TODO: Define the technical terminology, symbols, quantities and
% classifications used in this Deliverable.

\subsection{Literature Evidence}
\label{sec:res-ml-int-literature-evidence}

% TODO: Record evidence from peer-reviewed papers, authoritative datasets,
% standards, technical reports and primary documentation.
%
% Do not create a detached literature-summary list. Explain how each source
% supports, contradicts or limits a project decision.

\subsection{Governing Relations and Quantitative Definitions}
\label{sec:res-ml-int-governing-relations}

% TODO: Record relevant equations, dimensionless groups, empirical models,
% extraction rules or quantitative definitions.
%
% Use align environments for displayed equations.
% Every displayed equation must have a unique \label{eq:...}.
% Do not place punctuation after displayed equations.

\subsection{Machine-Learning Role}
\label{sec:res-ml-int-machine-learning-role}

% TODO: Complete this RES-ML Domain-specific subsection.

\subsection{Non-ML Baseline}
\label{sec:res-ml-int-non-ml-baseline}

% TODO: Complete this RES-ML Domain-specific subsection.

\subsection{Input and Output Representation}
\label{sec:res-ml-int-input-output-representation}

% TODO: Complete this RES-ML Domain-specific subsection.

\subsection{Dataset Requirements}
\label{sec:res-ml-int-dataset-requirements}

% TODO: Complete this RES-ML Domain-specific subsection.

\subsection{Model or Architecture Candidates}
\label{sec:res-ml-int-model-architecture-candidates}

% TODO: Complete this RES-ML Domain-specific subsection.

\subsection{Training and Validation Method}
\label{sec:res-ml-int-training-validation-method}

% TODO: Complete this RES-ML Domain-specific subsection.

\subsection{Evaluation Metrics}
\label{sec:res-ml-int-evaluation-metrics}

% TODO: Complete this RES-ML Domain-specific subsection.

\subsection{Generalisation and Uncertainty}
\label{sec:res-ml-int-generalisation-uncertainty}

% TODO: Complete this RES-ML Domain-specific subsection.

\subsection{Simulation and Optimisation Integration}
\label{sec:res-ml-int-simulation-optimisation-integration}

The deferred integration sequence is:
\begin{enumerate}
  \item consume validated Tohoku corridor and barrier-class outputs;
  \item specify a future candidate configuration or intervention;
  \item verify that the query lies inside the future surrogate domain;
  \item predict response and uncertainty;
  \item screen or rank reinforcement options;
  \item require direct hydrodynamic, structural or FSI simulation when
    uncertain or out of domain;
  \item add accepted verified simulations to the controlled dataset.
\end{enumerate}

\subsection{Candidate Approaches}
\label{sec:res-ml-int-candidate-approaches}

% TODO: Define the credible candidate approaches considered.

\subsection{Critical Comparison}
\label{sec:res-ml-int-critical-comparison}

% TODO: Compare candidates using physically and project-relevant criteria.
% Do not select an approach solely on popularity or implementation ease.

\subsection{Assumptions and Applicability}
\label{sec:res-ml-int-assumptions}

% TODO: State assumptions and the conditions under which they are valid.

\subsection{Limitations and Failure Conditions}
\label{sec:res-ml-int-limitations}

% TODO: State where the evidence, model or selected approach ceases to be
% reliable or applicable.

\subsection{Project-Specific Conclusions}
\label{sec:res-ml-int-conclusions}

Integration remains a governed future screening workflow. For the
current phase, accepted physical evidence flows from validated
simulation into structured outputs only; surrogate predictions do not
route decisions.

\subsection{Selected Approach and Rationale}
\label{sec:res-ml-int-selected-approach}

% TODO: Record the selected approach only after sufficient evidence exists.
% State the alternatives rejected and the reasons for rejection.

\subsection{Unresolved Decisions}
\label{sec:res-ml-int-unresolved-decisions}

% TODO: Record unresolved questions, required evidence and decision owners.

\subsection{Requirements and Handoffs}
\label{sec:res-ml-int-requirements}

% TODO: State requirements passed to other Research Domains, Software,
% verification activities or later execution work.

\subsection{Deliverable Completion Record}
\label{sec:res-ml-int-completion-record}

% TODO: Record whether the Deliverable is:
%
% - Not started
% - In progress
% - Draft complete
% - Reviewed
% - Accepted
%
% Acceptance requires:
%
% - the research questions to be answered;
% - claims to be supported by appropriate evidence;
% - assumptions and limitations to be explicit;
% - the project-specific conclusion to be stated;
% - downstream requirements to be traceable;
% - unresolved decisions to be closed or formally deferred.
